%%%%%%%% bonsai: MoE Minitron for TypeScript specialization %%%%%%%%%%%%%%%%%

\documentclass{article}

\usepackage{microtype}
\usepackage{graphicx}
\usepackage{subfigure}
\usepackage{booktabs}
\usepackage{hyperref}

\newcommand{\theHalgorithm}{\arabic{algorithm}}

\usepackage[accepted]{icml2025}

\usepackage{amsmath}
\usepackage{amssymb}
\usepackage{mathtools}
\usepackage{amsthm}
\usepackage[capitalize,noabbrev]{cleveref}

\theoremstyle{plain}
\newtheorem{theorem}{Theorem}[section]
\newtheorem{proposition}[theorem]{Proposition}
\newtheorem{lemma}[theorem]{Lemma}
\newtheorem{corollary}[theorem]{Corollary}
\theoremstyle{definition}
\newtheorem{definition}[theorem]{Definition}
\newtheorem{assumption}[theorem]{Assumption}
\theoremstyle{remark}
\newtheorem{remark}[theorem]{Remark}

\icmltitlerunning{bonsai}

\begin{document}

\twocolumn[
\icmltitle{Pruning and Distilling a Mixture-of-Experts Language Model \\
           for TypeScript Specialization}

\icmlsetsymbol{equal}{*}

\begin{icmlauthorlist}
\icmlauthor{Sartaj Bhuvaji}{su}
\end{icmlauthorlist}

\icmlaffiliation{su}{Seattle University, Seattle, WA, USA}

\icmlcorrespondingauthor{Sartaj Bhuvaji}{sbhuvaji@seattleu.edu}

\icmlkeywords{Mixture of Experts, Structured Pruning, Knowledge Distillation, Language Models, Code Generation}

\vskip 0.3in
]

\printAffiliationsAndNotice{}

\begin{abstract}
We apply the Minitron structured pruning and knowledge distillation recipe
\cite{minitron2024,minitron2024practice} to a Mixture-of-Experts (MoE)
language model, Qwen3-30B-A3B-Base, and specialize the result toward
TypeScript. Both Minitron papers prune only dense transformers; we extend
their forward-pass-only, activation-based importance-estimation philosophy
to a pruning axis neither paper covers -- expert pruning -- and report early
characterization results on the unpruned teacher model prior to any pruning
being performed. We find that a gate-weighted expert-output-norm criterion
(independently converged on by REAP \cite{reap2025} and NVIDIA's own
Minitron-MoE extension) correctly identifies previously-reported ``Super
Experts'' \cite{superexperts2025} as the top-ranked expert in their
respective layers, and that a cheap, single-forward-pass depth-pruning proxy
(Block Importance) correlates strongly (Pearson $r=0.91$) with the expensive,
$49$-forward-pass perplexity-based criterion from the original Minitron
paper, agreeing exactly on the three most critical layers. This is a living
document, updated as pruning and distillation experiments complete; no layer,
head, or expert has been pruned as of this writing.
\end{abstract}

\section{Introduction}
\label{sec:intro}

Minitron \cite{minitron2024} demonstrated that a large language model can be
compressed into a family of smaller models via structured pruning followed
by knowledge distillation, using up to $40\times$ fewer training tokens than
training each size from scratch while matching or exceeding the accuracy of
from-scratch baselines. A follow-up paper \cite{minitron2024practice} applied
the same recipe under a more realistic constraint -- no access to the
original pretraining data -- and introduced teacher correction and a
downstream-task-based depth-pruning criterion.

Both papers prune exclusively dense transformers. Every weight in a dense
model's MLP is applied to every token that passes through it. This is not
true of a Mixture-of-Experts (MoE) model: an MoE layer routes each token to a
small subset of experts (commonly 8 of 128, in the model we study here), so
most of the model's parameters are never touched by any given token. This
work asks what happens when the Minitron recipe is applied to a model with
this property, and in particular, how to score the importance of an axis --
expert pruning -- that has no precedent in either source paper.

Our target model is Qwen3-30B-A3B-Base
\cite{qwen3moe}: 48 decoder layers, 32 query / 4 key-value attention heads
(GQA), 128 routed experts per layer with 8 active per token, no shared
expert, approximately 30.5B total parameters of which approximately 3.4B
are active for any given token. We specialize this model toward TypeScript,
using \texttt{bigcode/the-stack-dedup}'s TypeScript split
\cite{stackdedup} for both calibration and retraining data, drawn from
disjoint parquet shards so the two pools never overlap.

This paper is written incrementally, alongside the project's own commit
history: each section is updated with real, measured results only, once the
corresponding experiment has actually run. As of this writing, no pruning
has occurred; every number reported below characterizes the unpruned
teacher model.

\section{Related Work}
\label{sec:related}

\textbf{Minitron.} \citet{minitron2024} prune Nemotron-4 15B to 8B and 4B
using purely activation-based importance estimation (no gradients), scoring
attention heads, MLP neurons, and embedding channels via a shared formula
that aggregates activations over a small (1024-sample) calibration set, and
scoring depth separately via perplexity delta or Block Importance (cosine
distance between a layer's input and output). \citet{minitron2024practice}
extend this to settings without the original pretraining data, introducing
teacher correction (lightly fine-tuning the teacher on the actual
distillation data before pruning) and a downstream-task-based depth-pruning
criterion using contiguous-block removal, evaluated on Llama 3.1 8B and
Mistral NeMo 12B.

\textbf{Expert pruning in MoE models.} Neither Minitron paper addresses
Mixture-of-Experts models. REAP \cite{reap2025} scores expert importance as
a gate-weighted expert-output norm, averaged over each expert's routed
tokens, and evaluates directly on Qwen3-30B-A3B (the model we study). NVIDIA's
Model-Optimizer library independently implements the same formula as part of
its own Minitron-MoE extension, used to compress Nemotron-3-Nano-30B-A3B.
Prior work on Mixtral-scale models found selection-frequency-based criteria
underperform output-magnitude-based ones \cite{naee2024}, consistent with
the choice we adopt (\cref{sec:experts}).

\textbf{Super Experts.} \citet{superexperts2025} report that Qwen3-30B-A3B
contains a small number of experts (three in the first four layers of the
model, by their numbering) whose rare but extremely large-magnitude outputs
anchor attention sinks; ablating them alone raises WikiText-2 perplexity
from 8.70 to 59.86. We treat this as a regression test for our own expert
importance metric (\cref{sec:experts}).

\section{Target Model and Data}
\label{sec:model-data}

\subsection{Model}
Qwen3-30B-A3B-Base~\cite{qwen3moe} is a base (not instruction-tuned) causal
language model with 48 \texttt{Qwen3MoeDecoderLayer} blocks, hidden size
2048, 32 query heads and 4 key-value heads (grouped-query attention,
head\_dim 128), and 128 routed experts per layer (\texttt{Qwen3MoeExperts},
moe\_intermediate\_size 768) with the top 8 active per token, no shared
expert. Every layer is sparse (\texttt{decoder\_sparse\_step}=1,
\texttt{mlp\_only\_layers} empty). Experts within a layer are not separate
submodules but two batched parameter tensors
(\texttt{gate\_up\_proj}, \texttt{down\_proj}), with the expert index as the
leading dimension -- pruning an expert means slicing that dimension.

\subsection{Data}
Calibration and retraining data are both drawn from
\texttt{bigcode/the-stack-dedup}'s \texttt{data/typescript} split (71
parquet shards). The first shard is reserved for calibration, the remaining
70 for retraining; because the split is shard-level rather than row-level,
the two pools are disjoint by construction, which we additionally verify
empirically against each sample's \texttt{hexsha}.

\section{Baseline Evaluation of the Unpruned Teacher}
\label{sec:baseline}

Before any pruning decision, we establish the unpruned model's performance
on the evaluation suite every subsequent pruning candidate will be measured
against. Following both Minitron papers' evaluation settings (temperature
0.2, nucleus sampling top-$p=0.95$ for generation tasks), we report:

\begin{table}[t]
\caption{Baseline evaluation of Qwen3-30B-A3B-Base, unpruned.}
\label{tab:baseline}
\centering
\small
\begin{tabular}{lcc}
\toprule
Task & Setting & Score \\
\midrule
HumanEval-TS & pass@1 & 0.6708 \\
MBPP-TS & pass@1 & 0.7333 \\
MMLU & 5-shot, acc & 0.8101 \\
ARC-Challenge & 25-shot, acc & 0.6536 \\
ARC-Challenge & 25-shot, acc\_norm & 0.6749 \\
Winogrande & 5-shot, acc & 0.7798 \\
\bottomrule
\end{tabular}
\end{table}

HumanEval-TS and MBPP-TS (both MultiPL-E TypeScript translations) are
treated as the primary, gating tier, since they measure the capability this
project specializes toward. MMLU, ARC-Challenge, and Winogrande are tracked
but not optimized for, giving a comparison point against both Minitron
papers' own general-knowledge tables without over-indexing on general
capability at the expense of TypeScript performance.

\section{Importance Estimation}
\label{sec:importance}

Both Minitron papers score importance from forward-pass activations only,
never gradients, over a small calibration set. We follow the same
philosophy for every axis, including the one neither paper covers.

\subsection{Attention Heads}
Head importance uses the original formula unchanged: the per-token,
per-head output vector's $L_2$ norm, aggregated over every token in the
calibration set. Across all 48 layers (48 calibration samples, mean length
177 tokens), scores range from 2.39 to 967.72, with the single largest value
at the model's final layer (layer 47).

\subsection{Experts}
\label{sec:experts}
Neither Minitron paper offers a formula for this axis. We adopt
$S_j = \text{mean}_{x \in X_j}\left[ g_j(x) \cdot \lVert f_j(x) \rVert_2
\right]$, where $X_j$ is the set of tokens routed to expert $j$, $g_j(x)$ is
that token's renormalized routing weight, and $f_j(x)$ is the expert's
output -- the formula independently reached by both REAP \cite{reap2025} and
NVIDIA Model-Optimizer's Minitron-MoE extension, in preference to two
alternatives considered (selection frequency alone; router logit magnitude
alone), both reported to underperform output-magnitude-based scoring in
prior work \cite{naee2024}.

As a validation check, not a designed experiment: all three Super Experts
reported by \citet{superexperts2025} for this exact model (layer 1 / expert
68, layer 2 / expert 92, layer 3 / expert 82) rank \emph{first} among all
128 experts in their respective layers under this metric, on real
calibration data, at two different sample sizes (24 and 48 samples). On
average across all 48 layers, the top $\sim$16\% of experts (20 of 128)
already account for roughly a third of a layer's total importance mass,
indicating substantial redundancy among the remaining experts as a
candidate for pruning.

\subsection{Depth: Block Importance and Perplexity}
\label{sec:depth}
We implement both depth-scoring criteria from \citet{minitron2024}: Block
Importance (BI), $1 - \text{cosine\_similarity}(\text{layer\_input},
\text{layer\_output})$, computable in the same forward pass that produces
the two activations; and a perplexity-delta criterion, which removes one
layer at a time (via a context manager that temporarily makes the layer act
as identity, restoring it afterward) and measures the resulting increase in
held-out cross-entropy loss relative to the unmodified baseline. The latter
costs $L+1$ forward passes for an $L$-layer model, versus BI's one, and is
run on a smaller calibration sample (16 samples, vs.\ 48 for BI) for that
reason.

Both criteria independently identify the same pattern: the first and last
layers of the model matter far more than the middle two-thirds. BI's
highest score is at layer 0 (0.488) and second-highest at layer 47 (0.415);
its lowest is layer 12 (0.011). The perplexity criterion shows the same
edges-matter-most shape at a far more extreme scale -- layer 0's removal
costs 1.99 nats of cross-entropy, next-highest is layer 47 at 0.60 nats,
and the lowest (layer 14) is statistically indistinguishable from zero
($-0.0003$ nats).

Across all 48 layers, the Pearson correlation between BI and the perplexity
delta is $r=0.91$, and the two criteria pick the identical top-3 most
critical layers (0, 47, 1). Given the 49$\times$ cost difference between the
two, this agreement is the paper's most directly actionable early finding:
Block Importance can be trusted as a cheap stand-in for perplexity-based
scoring when iterating over many candidates later in this project, reserving
the expensive criterion for confirming a final choice rather than for every
candidate evaluated along the way.

\section{Discussion and Open Questions}
\label{sec:discussion}

Several design decisions remain open and are logged here rather than
resolved prematurely:

\begin{itemize}
\item \textbf{Downstream-depth-pruning proxy.} \citet{minitron2024practice}'s
  preferred depth criterion ranks contiguous layer blocks by downstream-task
  accuracy after removal (Winogrande, in their case), not by BI or
  perplexity. No TypeScript equivalent of Winogrande exists; HumanEval-TS or
  MBPP-TS pass-rate on partially-pruned models is the direct analogue but
  expensive per candidate block, and a cheaper token-level-accuracy proxy is
  under consideration.
\item \textbf{Uniform vs.\ per-layer expert counts.} We currently plan
  uniform expert counts per layer for the first pruning pass, following
  published evidence that non-uniform allocation did not outperform uniform
  allocation on this specific model in at least one independent evaluation,
  though this has not yet been tested within this project directly.
\item \textbf{Interaction between axes.} \citet{minitron2024} find that
  depth and width importance rankings can flip after retraining relative to
  their zero-shot rankings; whether the same holds when experts are combined
  with depth and width has no precedent to draw on and will be tested
  empirically once single-axis pruning is complete.
\end{itemize}

No results are reported for these open questions; they are listed to make
the project's actual state legible, not to imply an answer.

\section*{Acknowledgements}
This is an independent research project with no external funding.

\bibliography{references}
\bibliographystyle{icml2025}

\end{document}
